\chapter{Contexte du stage et présentation du projet}

\section{Présentation de l'organisme d'accueil}

\noindent
Le stage a été réalisé au sein de \entreprise{}. Les informations détaillées concernant l'entreprise, son historique, son organisation interne, ses activités principales et ses références clients seront complétées après validation des données officielles par l'organisme d'accueil.

\noindent
Dans le cadre de ce rapport, l'entreprise est considérée comme un environnement professionnel orienté vers la conception et le développement de solutions web. Le projet \projectName{} s'inscrit dans une logique de digitalisation d'un processus métier concret : le recrutement dans le secteur CHR.

\section{Contexte métier}

\noindent
Le secteur CHR, qui regroupe les cafés, hôtels et restaurants, représente un domaine où les besoins en recrutement sont fréquents. Les établissements recherchent souvent des profils opérationnels tels que serveurs, cuisiniers, réceptionnistes, femmes de chambre, baristas ou commis de cuisine.

\noindent
Cependant, le processus de recrutement reste parfois dispersé. Les recruteurs publient des annonces sur plusieurs canaux, reçoivent des CV de manière non centralisée et rencontrent des difficultés à qualifier rapidement les candidats. De leur côté, les candidats manquent d'un espace simple pour présenter leur profil, déposer leur CV et suivre leurs candidatures.

\section{Problématique}

\noindent
La problématique principale du projet peut être formulée comme suit :

\begin{quote}
Comment concevoir une plateforme web simple, structurée et adaptée au secteur CHR, permettant de mettre en relation les candidats et les recruteurs tout en améliorant le suivi des candidatures et la qualification des profils ?
\end{quote}

\noindent
Cette problématique implique plusieurs besoins : centraliser les offres, rendre le profil candidat réutilisable, éviter la répétition de l'upload du CV, permettre au recruteur d'évaluer les profils, intégrer un quiz métier et offrir à l'administrateur un outil de modération.

\section{Objectifs du projet}

\noindent
Les objectifs de \projectName{} sont les suivants :

\begin{itemize}
    \item créer une plateforme web de recrutement dédiée au secteur CHR ;
    \item permettre aux candidats de compléter un profil professionnel et d'ajouter un CV ;
    \item permettre aux recruteurs de publier et gérer leurs offres ;
    \item intégrer un système de candidatures avec statut et score de quiz ;
    \item proposer un espace administrateur pour la modération des offres ;
    \item améliorer l'expérience utilisateur à travers une interface claire et responsive.
\end{itemize}

\section{Périmètre fonctionnel}

\noindent
Le projet couvre trois espaces principaux :

\begin{table}[h!]
\centering
\renewcommand{\arraystretch}{1.4}
\begin{tabular}{|p{3cm}|p{10cm}|}
\hline
\textbf{Espace} & \textbf{Fonctionnalités principales} \\
\hline
Candidat & Profil professionnel, CV, photo optionnelle, consultation des offres, postulation, quiz, suivi des candidatures, favoris et notifications. \\
\hline
Recruteur & Profil établissement, création d'offres, image optionnelle, quiz optionnel, consultation des candidatures, CV, acceptation/refus et quota premium. \\
\hline
Administrateur & Statistiques globales, liste des utilisateurs, modération des offres, suspension avec motif et réactivation. \\
\hline
\end{tabular}
\caption{Périmètre fonctionnel de SmartJobs}
\label{tab:perimetre_fonctionnel}
\end{table}

\section{Contraintes du projet}

\noindent
Le projet devait respecter plusieurs contraintes :

\begin{itemize}
    \item conserver une architecture web claire entre frontend et backend ;
    \item sécuriser les accès selon les rôles ;
    \item limiter les uploads de fichiers aux formats autorisés ;
    \item éviter les candidatures dupliquées ;
    \item rendre l'interface utilisable en mode sombre et en mode clair ;
    \item préparer une démonstration stable pour la soutenance.
\end{itemize}
